\section{Coherent packets, joint tensor layers, and their exact spectral cost}
\label{sec:coherent-tensor-integration}

This section integrates the complete arithmetic interpolation construction
of the 12 September coherent-interpolation delivery, whose source revision
is \nolinkurl{3bff1b2bd87dfa2592f65c5ccc231f6ee9541dfc}, and proves its
comparison with the preceding kernel and departure calculations.  All
zeros in a packet retain their complete multiplicities.  A reflected
packet is required only in statements explicitly using its reflection.
The original spaces \(V,\mathscr B\), map \(\Theta\), generator
\(D=-x\partial_x\), and Hilbert observation \(L^2(dx)\) remain fixed.

\subsection{The arithmetic unit and the coherent one-variable maps}

For \(h(s)=\prod_{\rho\in Z}(s-\rho)^{m_\rho}\), write
\[
 d=\deg h,\quad E_h=\mathbb C[s]/h,\quad A_h=M_s,\quad
 v_h=g/h,\quad \upsilon_h=j_hv_h,\quad
 \varepsilon_h=\upsilon_h^{-1},\qquad g=2\xi.
 \tag{CT.1}
\]
Complete multiplicities make \(\upsilon_h\) invertible.  Let
\(F_h\in\mathscr B\) be the original inverse satisfying
\(h(D)F_h=f_0\) and \(\mathcal MF_h=g/h\), and put
\(\mathcal T_hP=P(D)F_h\).  The integral construction of
\(F_h\), its estimates at both ends, and injectivity of the Mellin
transform are proved in Section~\ref{sec:arithmetic-input}.  They give
\[
 \mathcal T_h(hP)=\Theta(P(D)\phi_*),\qquad
 J_h\mathcal T_hP=\upsilon_h[P]_h.                       \tag{CT.2}
\]
Let \(C_N:\mathcal P_N\to E_h\) be the full remainder map in
the original monomial coordinates, let \(M_{h,N}\) be the Gram
matrix of \(\mathcal P_N\) for
\(d\nu_h=|g(1/2+it)/h(1/2+it)|^2dt/(2\pi)\), and use
\(U_a\) for multiplication by \(a\) on \(E_h\).  The two
kernel matrices occurring in the constructions have the exact types
\[
 K_N=C_NM_{h,N}^{-1}C_N^*,\qquad
 J_{h,N}=U_{\upsilon_h}C_N,\qquad
 \mathcal C_{h,N}=U_{\upsilon_h}K_NU_{\upsilon_h}^*.
 \tag{CT.3}
\]
Consequently their actual representative and metric coincide:
\[
 \begin{split}
 M_{h,N}^{-1}J_{h,N}^*\mathcal C_{h,N}^{-1}
   &=M_{h,N}^{-1}C_N^*K_N^{-1}U_{\varepsilon_h},\\
 G_{h,n}=\mathcal C_{h,n+d}^{-1}
   &=U_{\varepsilon_h}^*K_{n+d}^{-1}U_{\varepsilon_h}.
 \end{split}                                                   \tag{CT.4}
\]
Indeed invert the congruence in (CT.3) and cancel the adjacent
inverse factors.  The first line is exactly the polynomial numerator
in (TB.7)--(TB.9).  Its difference from
\(\operatorname{rem}_h(\varepsilon_hu)\) is divisible by \(h\)
and has quotient of degree at most \(n\).  Applying (CT.2) gives
the same original theta correction.  Thus the arithmetic unit remains
in every coefficient map and every source representative.

For completeness, the minimum assertion follows directly without an
existence assumption.  For any positive Gram matrix \(M\) and
surjective \(J\), set
\(G=(JM^{-1}J^*)^{-1}\) and \(L=M^{-1}J^*G\).  Then
\(JL=I\), \(L^*ML=G\), and for \(Jz=0\),
\(z^*ML=z^*J^*G=0\).  Every feasible vector is \(Lu+z\), so
\[
 (Lu+z)^*M(Lu+z)=u^*Gu+z^*Mz.                            \tag{CT.5}
\]
Positivity makes this the unique minimum.  In (CT.3), positivity
holds because a nonzero polynomial cannot vanish almost everywhere
on the original integration line; surjectivity follows from complete
polynomial remainders.  The kernel is exactly
\(h\mathcal P_n\) when \(N=n+d\).  Hence
\(R_{h,n}=(I-P_{L_n})s_h\), where
\(s_h=\mathcal T_h\operatorname{rem}_h(\varepsilon_h\,\cdot)\)
and \(L_n=\Theta\operatorname{span}(\phi_*,\ldots,D^n\phi_*)\).

If \(h\mid H\) are full packets, put \(m=H/h\).  The factors
\(h,m\) are coprime.  The original CRT injection
\(\iota:E_h\to E_H\) inserts zero jets at the additional
centres, and satisfies \(A_H\iota=\iota A_h\).
Multiplication \(P\mapsto mP\) is an isometry of the declared
polynomial images, since \((g/H)mP=(g/h)P\); it satisfies
\[
 \mathcal T_H(mP)=\mathcal T_hP,\qquad
 J_H\mathcal T_H(mP)=\iota J_h\mathcal T_hP.              \tag{CT.6}
\]
At the added centres the second identity follows from the full zeros
of \(g/h\); at the old centres it is the identity of entire
functions.  Both numerators for
\(s_H\iota=\mathcal T_H(m\operatorname{rem}_h
(\varepsilon_h\,\cdot))\) have degree below \(\deg H\)
and the same complete jets, proving equality.  Since \(L_n\) is
independent of the packet,
\[
 R_{H,n}\iota=R_{h,n},\quad
 \iota^*G_{H,n}\iota=G_{h,n},\quad
 \iota^*W_{H,n}\iota=W_{h,n}.                            \tag{CT.7}
\]
These are congruences, with no interchange of inversion and
compression.  Their split maps at a retained support \(\lambda\)
are \((\lambda,u)\mapsto(\lambda,\iota u)\); all added zero
jets are represented zeros at that support.

\subsection{The original reflection and the exact diagonal coefficient}

The original arithmetic density \(|g(1/2+it)|^2\) is even in
\(t\), since the real Taylor coefficients of \(g\) give
\(g(1/2-it)=\overline{g(1/2+it)}\).  If \(Q_j\) is the
monic orthogonal polynomial for that density in the variable \(s\),
the polynomial \((-1)^jQ_j(1-s)\) is monic and orthogonal to
all lower degrees, by the change \(t\mapsto-t\).  Uniqueness
therefore gives
\[
 Q_j(1-s)=(-1)^jQ_j(s),\qquad
 \int t\,|Q_j(1/2+it)|^2|g(1/2+it)|^2\frac{dt}{2\pi}=0.
 \tag{CT.8}
\]
The diagonal recurrence coefficient is thus exactly
\(\langle Q_j,sQ_j\rangle/\|Q_j\|^2=1/2\).
For lower degrees the identity \(D^*+D=I\) eliminates every
coefficient except the adjacent one and gives its sign and actual
norm ratio.  More explicitly, \(sQ_j\) is orthogonal to
\(Q_i\) for \(i\le j-2\), because its inner product is the
inner product of \(Q_j\) with \((1-s)Q_i\).  For \(i=j-1\)
the leading term gives coefficient \(-\|Q_j\|^2/\|Q_{j-1}\|^2\).
This proves, with the original vectors and norms,
\[
 D u_j=u_{j+1}+\tfrac12u_j
       -\frac{\mathfrak h_j}{\mathfrak h_{j-1}}u_{j-1},
 \qquad u_j=Q_j(D)f_0,\quad \mathfrak h_j=\|u_j\|^2.
 \tag{CT.9}
\]
The final term is omitted at \(j=0\).  The even density in
(CT.8) is the additional reason the imaginary part of the diagonal
coefficient vanishes; the adjoint identity by itself fixes only its
real part.

For a packet stable under \(\rho\mapsto1-\overline\rho\),
including the same complete orders, the original maps are
\[
 \mathfrak j_hu(s)=\overline{u(1-\overline s)}\pmod h,\qquad
 \mathfrak j_{\mathscr B}F(x)=x^{-1}\overline{F(1/x)}.
\]
The second is antiunitary for \(dx\), and
\(D\mathfrak j_{\mathscr B}=\mathfrak j_{\mathscr B}(I-D)\).  The identities
\(h^\dagger=(-1)^dh\) and \(g^\dagger=g\) imply
\((g/h)^\dagger=(-1)^d g/h\); they give
\(\mathfrak j_{\mathscr B}s_h=s_h\mathfrak j_h\) by the same complete jets and
degree bounds as above.  Equation (CT.9), or invariance of the
polynomial space under \(s\mapsto1-s\), shows
\(\mathfrak j_{\mathscr B}L_n=L_n\).  The unique minimum (CT.5) then
commutes with these maps.  If \(\mathfrak j_hu=C_h\overline u\),
\[
 C_h^*G_{h,n}C_h=\overline{G_{h,n}},\quad
 A_hC_h=C_h(I-\overline{A_h}),\quad
 C_h^*W_{h,n}C_h=-\overline{W_{h,n}}.                     \tag{CT.10}
\]
Substitution into a quadratic form proves that an upper bound and its
reflected lower bound use this same \(G_{h,n}\).

\subsection{The joint total-degree source and its relation layer}

Fix \(k\ge1\), \(M\ge k(d-1)\), and set
\[
 E^{(k)}=E_h^{\otimes k},\quad A^{(k)}=\sum_{j=1}^kA_{h,j},
 \quad \mathcal P_M^{(k)}=\mathbb C[s_1,\ldots,s_k]_{\le M},
 \quad \mathcal T_h^{(k)}P=P(D_1,\ldots,D_k)F_h^{\otimes k}.
\]
Here \(\le M\) means total degree.  The full jet map is
reduction by each \(h(s_j)\), followed by multiplication by
\(\upsilon_h^{\otimes k}\), and its kernel is
\[
 \mathcal I_{h,k,M}=\mathcal P_M^{(k)}
            \cap(h(s_1),\ldots,h(s_k)).                 \tag{CT.11}
\]
Successive monic division in the declared order
\(s_1,\ldots,s_k\) gives
\(P=\sum_jh(s_j)P_j\) in this kernel, with
\(\deg P_j\le M-d\).  Division lowers total degree at every
replacement, so this bound holds at every division step.  In the
original complex \([V\to\mathscr B]\), placed in degrees
\(0,1\), the primitive for the \(j\)-th term uses
\((-1)^{j-1}\phi_*\) in the \(j\)-th factor and \(F_h\)
in the others, with the polynomial \(P_j(D_1,\ldots,D_k)\)
applied to the whole tensor.  The preceding \(j-1\) factors have
degree one, so the tensor differential supplies the further factor
\((-1)^{j-1}\).  Its boundary is exactly
\(+\mathcal T_h^{(k)}(h(s_j)P_j)\).

Write \(\mathcal B_M=\mathcal T_h^{(k)}\mathcal I_{h,k,M}\)
inside \(\mathcal H_k=L^2(d^kx)\).  The arithmetic moment
matrix in the retained monomials is
\[
 (M_{k,M})_{\alpha\beta}=
   \prod_{j=1}^k\int\overline{s^{\alpha_j}}s^{\beta_j}
                   \,d\nu_h,\quad
 G_M=(J_{k,M}M_{k,M}^{-1}J_{k,M}^*)^{-1}.                 \tag{CT.12}
\]
Product integration proves the first formula.  A polynomial vanishing
almost everywhere for this product measure is zero, by successive
one-variable uniqueness, proving positivity.  The complete remainder
basis has total degree at most \(k(d-1)\), proving surjectivity.
Thus (CT.5) proves
\[
 R_M=(I-P_{\mathcal B_M})s_h^{\otimes k},\qquad
 J^{(k)}R_M=I,\qquad q^{(k)}R_M=\sigma_h^{\otimes k}.
 \tag{CT.13}
\]
The kernel dimension is \(\binom{M+k}{k}-d^k\).  Since
\(\mathcal T_h^{(k)}\) is injective by its entire Mellin
multiplier, the orthogonal relation layer
\(\mathcal E_M=\mathcal B_{M+1}\cap\mathcal B_M^\perp\)
has exactly
\[
 \dim\mathcal E_M=\binom{M+k}{k-1}.                      \tag{CT.14}
\]
Define all of the following maps on the original \(E^{(k)}\):
\[
 B_M=D^{(k)}R_M-R_MA^{(k)},\quad
 Y_M=P_{\mathcal E_M}R_M,\quad
 C_M=(I-P_{\mathcal B_M})B_M.                            \tag{CT.15}
\]
The complete jet of \(B_M\) is zero and its numerator has
degree at most \(M+1\), so
\(B_M(E^{(k)})\subseteq\mathcal B_{M+1}\).
Both \(Y_M,C_M\) therefore have target \(\mathcal E_M\).
Integration by parts on these rapidly decreasing functions gives
\((D^{(k)})^*+D^{(k)}=kI\).  Substitution into
\(R_M^*(D^{(k)}+(D^{(k)})^*)R_M=kG_M\), followed by
orthogonality to \(\mathcal B_M\), proves
\[
 \begin{split}
 W_M&=(A^{(k)})^*G_M+G_MA^{(k)}-kG_M
       =-(Y_M^*C_M+C_M^*Y_M),\\
 \Delta_M&=G_M-G_{M+1}=Y_M^*Y_M\succeq0.
 \end{split}                                                   \tag{CT.16}
\]
For the second identity use the orthogonal decomposition of
\(\mathcal B_{M+1}\) and \(R_{M+1}=R_M-Y_M\).
The actual quotient map
\([z]\mapsto(I-P_{\mathcal B_M})z\) identifies
\(\mathcal B_{M+1}/\mathcal B_M\) with \(\mathcal E_M\),
with inverse the inclusion followed by quotient.  A transported
boundary is zero at its next support and remains represented there;
this map sends no supported input to external \(\tau\).

For a reflection-stable packet, put \(C_{h,k}=C_h^{\otimes k}\).
The original product reflection acts on a numerator by
\((-1)^{kd}\overline{P(1-\overline{s_1},\ldots,
1-\overline{s_k})}\).  This retains total degree and preserves
the ideal in (CT.11), with its displayed factor.  Its action on
\(\mathcal H_k\) is antiunitary.  The unique minimum (CT.5)
therefore commutes with reflection and gives
\[
 C_{h,k}^*G_MC_{h,k}=\overline{G_M},\quad
 A^{(k)}C_{h,k}=C_{h,k}(kI-\overline{A^{(k)}}),\quad
 C_{h,k}^*W_MC_{h,k}=-\overline{W_M}.                    \tag{CT.16a}
\]
If \(W_Mu=\lambda G_Mu\), conjugation followed by
\((C_{h,k}^*)^{-1}\) in the last identity gives
\(W_MC_{h,k}\overline u=-\lambda G_MC_{h,k}\overline u\).
All generalized eigenvalues are real because \(G_M>0\).
This invertible antilinear map pairs their full multiplicities.
The source-layer maps themselves retain the reflection types.  If
\(\mathfrak j_E=\mathfrak j_h^{\otimes k}\) and
\(\mathfrak j_{\mathcal E}\) is the restriction of the
original function reflection to \(\mathcal E_M\), then
\[
 \mathfrak j_{\mathcal E}Y_M=Y_M\mathfrak j_E,
 \qquad \mathfrak j_{\mathcal E}C_M=-C_M\mathfrak j_E.
 \tag{CT.16b}
\]
Both boundary projections commute with reflection, proving the first
identity.  For the second use
\(D^{(k)}\mathfrak j_{\mathcal H}
=\mathfrak j_{\mathcal H}(kI-D^{(k)})\),
\(A^{(k)}\mathfrak j_E=\mathfrak j_E(kI-A^{(k)})\),
and \(\mathfrak j_{\mathcal H}R_M=R_M\mathfrak j_E\)
in (CT.15).  Expanding its two terms gives
\(\mathfrak j_{\mathcal H}B_M=-B_M\mathfrak j_E\);
the projection then proves (CT.16b).

\subsection{Nested packets on the joint source: the additional projection}

Let \(h\mid H\), \(q=\deg(H/h)\),
\(N=M+kq\), and \(\iota_k=\iota_{hH}^{\otimes k}\).
Multiplication of a numerator by
\(\prod_{j=1}^k(H/h)(s_j)\) gives a degree-preserving-with-shift
map \(\mathcal P_M^{(k)}\to\mathcal P_N^{(k)}\), with
\[
 \mathcal T_H^{(k)}\bigl(\textstyle\prod_j(H/h)(s_j)P\bigr)
       =\mathcal T_h^{(k)}P,\qquad
 J_H^{(k)}\mathcal T_h^{(k)}P
       =\iota_kJ_h^{(k)}\mathcal T_h^{(k)}P.              \tag{CT.17}
\]
The second equality follows on pure monomials from (CT.6) and then
by linearity.  In particular it maps the smaller relation space
into the larger one in the same \(\mathcal H_k\).
Let
\[
 \mathcal F=\mathcal B_{H,k,N}\cap\mathcal B_{h,k,M}^{\perp},
 \qquad Z=P_{\mathcal F}R_{h,k,M}.
\]
Since \(s_H^{\otimes k}\iota_k=s_h^{\otimes k}\), the exact
comparison is
\[
 \begin{split}
 R_{H,k,N}\iota_k&=R_{h,k,M}-Z,\\
 \iota_k^*G_{H,k,N}\iota_k&=G_{h,k,M}-Z^*Z,\\
 \iota_k^*W_{H,k,N}\iota_k
   &=W_{h,k,M}-\bigl((A_h^{(k)})^*Z^*Z
                         +Z^*ZA_h^{(k)}-kZ^*Z\bigr).
 \end{split}                                                   \tag{CT.18}
\]
Indeed nested orthogonal projections commute, and their difference
is \(P_{\mathcal F}\); apply them to the common section to prove
the first line.  Pythagoras proves the second; intertwining of
\(A^{(k)}\) and direct expansion prove the third.  Also
\[
 \dim\mathcal F=
 \binom{N+k}{k}-(\deg H)^k-\binom{M+k}{k}+d^k.           \tag{CT.19}
\]
This is the difference of the exact kernel dimensions.  For
\(k=1\), \(M=n+d\), both boundary spaces are the original
\(L_n\), so \(\mathcal F=0\) and (CT.18) recovers (CT.7).
For joint tensors (CT.18) retains every additional allowed boundary
and calculates its effect on the control form.

\subsection{Positive layer rank and the complete tensor spectral departure}

List \(\rho_1,\ldots,\rho_d\) with all algebraic
multiplicities, put \(\delta_i=\operatorname{Re}\rho_i-1/2\),
\(\sigma_1=\sum_i\delta_i\), \(\sigma_2=\sum_i\delta_i^2\),
and define
\[
 \begin{split}
 S_M&=G_M^{-1/2}W_MG_M^{-1/2},\quad
       \epsilon_M=\|S_M\|_{\rm op},\quad
       r_M=\operatorname{rank}\Delta_M,\\
 \mathfrak d_M^{(k)}&=
 \operatorname{Tr}(G_M^{-1}(A^{(k)})^*G_MA^{(k)})
    -\sum_{i_1,\ldots,i_k}\left|\sum_{j=1}^k\rho_{i_j}\right|^2,\\
 \mathcal D_k&=\sum_{i_1,\ldots,i_k}
                          \left|\sum_{j=1}^k\delta_{i_j}\right|.
 \end{split}                                                   \tag{CT.20}
\]
All sums are ordered and retain repeated entries.  The isometry
\(x\mapsto G_M^{1/2}x\) in this calculation leaves the original
jet coordinates, source functions, metric and layer maps unchanged.
For a map into \(\mathcal E_M\), the Frobenius symbol below
denotes its Hilbert--Schmidt norm for the layer's inherited
\(L^2(d^kx)\) inner product.

\begin{theorem}[Exact tensor departure and relative layer cost]
One has \(\mathfrak d_M^{(k)}\ge0\), and exactly
\[
 \begin{split}
 \operatorname{Tr}S_M&=2k d^{k-1}\sigma_1,\\
 \operatorname{Tr}(S_M^2)
    &=4\{kd^{k-1}\sigma_2+k(k-1)d^{k-2}\sigma_1^2\}
                        +2\mathfrak d_M^{(k)}.
 \end{split}                                             \tag{CT.21}
\]
The positive and negative inertias of \(W_M\) each obey
\[
 n_+(W_M),n_-(W_M)\le
 \min(\operatorname{rank}Y_M,\operatorname{rank}C_M),
 \quad \operatorname{rank}Y_M=r_M
           \le\binom{M+k}{k-1}.                         \tag{CT.22}
\]
The complete ordered spectral cost satisfies
\[
 \boxed{\mathcal D_k\le\tfrac12\|S_M\|_1
      \le \|Y_MG_M^{-1/2}\|_F\|C_MG_M^{-1/2}\|_F,}
 \qquad \mathcal D_k\le r_M\epsilon_M.                  \tag{CT.23}
\]
In particular, the actual source matrices satisfy
\[
 \mathcal D_k^2\le
 \operatorname{Tr}(G_M^{-1}\Delta_M)
                  \operatorname{Tr}(G_M^{-1}C_M^*C_M).
 \tag{CT.24}
\]
If the packet is reflection-stable, put
\(p_M=n_+(W_M)=n_-(W_M)\).  Then
\[
 p_M\le\min\{r_M,\operatorname{rank}C_M,
                         \lfloor d^k/2\rfloor\},\qquad
 \mathcal D_k\le p_M\epsilon_M.                         \tag{CT.24a}
\]
In this case (CT.21) becomes
\(\operatorname{Tr}(S_M^2)=4kd^{k-1}\sigma_2
+2\mathfrak d_M^{(k)}\).  When \(p_M>0\), it implies
\[
 \epsilon_M^2\ge
       \frac{2kd^{k-1}\sigma_2+\mathfrak d_M^{(k)}}{p_M}.
 \tag{CT.25}
\]
\end{theorem}

\begin{proof}
Triangularize \(A_h\) over \(\mathbb C\).  The tensor product
of this basis makes \(A^{(k)}\) triangular, with diagonal entries
\(\sum_j\rho_{i_j}\).  Thus these are all its eigenvalues
with algebraic multiplicity; its nilpotent entries have not been
discarded.  Now apply the unitary triangularization argument of
(MD5) to \(T=G_M^{1/2}A^{(k)}G_M^{-1/2}\).  The sum of
squared absolute values of its strictly upper triangular entries is
exactly \(\mathfrak d_M^{(k)}\), proving nonnegativity.
Since \(S_M=T+T^*-kI\), expansion of its squared trace gives
\[
 \operatorname{Tr}(S_M^2)=
    4\sum_{i_1,\ldots,i_k}(\sum_j\delta_{i_j})^2
                      +2\mathfrak d_M^{(k)}.
\]
The \(k\) diagonal terms in this ordered sum contribute
\(kd^{k-1}\sigma_2\); the \(k(k-1)\) ordered cross terms
contribute \(k(k-1)d^{k-2}\sigma_1^2\).  The linear trace is
computed in the same way.  This proves (CT.21), including \(k=1\).

Write \(\widehat Y=Y_MG_M^{-1/2}\) and
\(\widehat C=C_MG_M^{-1/2}\).  Then
\(S_M=-(\widehat Y^*\widehat C+
\widehat C^*\widehat Y)\).  Its form is zero on
\(\ker\widehat Y\).  A positive definite subspace of dimension
greater than \(\operatorname{rank}\widehat Y\) intersects this
kernel nontrivially, a contradiction; the same argument holds for
a negative definite subspace and with \(\widehat C\) in place
of \(\widehat Y\).  Sylvester inertia under the displayed
invertible congruence proves the first part of (CT.22).
The identity \(\Delta_M=Y_M^*Y_M\) gives equality of their
kernels, hence their ranks, proving the rest.

In the unitary triangular basis of \(T\), the diagonal of
\(S_M\) consists of \(2\sum_j\delta_{i_j}\).
For any Hermitian matrix \(S\), its spectral decomposition gives
\(|S_{ii}|\le\sum_a|\lambda_a||U_{ia}|^2\).  Sum over \(i\)
to obtain \(\sum_i|S_{ii}|\le\sum_a|\lambda_a|=\|S\|_1\).
This proves the first inequality in (CT.23).
For finite matrices \(B,C\) with a common row space,
\(\|B^*C\|_1\le\|B\|_F\|C\|_F\): choose a singular-value
decomposition of \(B^*C\), use its left and right singular
vectors to express its trace norm as
\(|\operatorname{Tr}(U B^*C)|\) for a unitary \(U\), and
apply Cauchy--Schwarz to the entries of \(BU^*\) and \(C\).
The triangle inequality now proves the second inequality of
(CT.23).  Finally (CT.22) bounds the number of nonzero eigenvalues
of \(S_M\) by \(2r_M\), so
\(\|S_M\|_1\le2r_M\epsilon_M\).  Squaring the middle
bound and writing its two Frobenius norms as traces proves (CT.24).

For a reflected packet \(\sigma_1=0\), with every multiplicity
paired.  The exact reflection (CT.16a) pairs the nonzero eigenvalues
of \(S_M\), giving equal positive and negative counts \(p_M\).
Equation (CT.22) proves the bounds on \(p_M\), and there are
exactly \(2p_M\) nonzero eigenvalues of \(S_M\), each of square
at most \(\epsilon_M^2\).  Their trace norm is at most
\(2p_M\epsilon_M\), proving (CT.24a); their squared trace
proves (CT.25).  If \(p_M=0\), the pairing implies
\(W_M=0\) directly, and (CT.21) states that all the
nonnegative terms vanish; no division by \(p_M\) is made.
\end{proof}

Every quantity in (CT.24) is attached to the same arithmetic moment
matrix and its next original relation layer.  In particular,
\(0\preceq\Delta_M\prec G_M\), because \(G_{M+1}>0\),
so all eigenvalues of \(G_M^{-1/2}\Delta_MG_M^{-1/2}\) lie
in \([0,1)\), and its trace is strictly smaller than \(r_M\)
when \(r_M>0\).  Its actual value is retained in (CT.24).
Equations (CT.21)--(CT.25) calculate necessary relative layer costs
and the full nilpotent departure for the constructed tensor family.
They assert no unproved upper bound on these arithmetic quantities.

There is an exact determinant refinement of the same source loss.
For \(r_M>0\), put
\[
 \pi_M=\frac{\det G_{M+1}}{\det G_M},\qquad
 \tau_M=\operatorname{Tr}(G_M^{-1}\Delta_M),\qquad
 \chi_M=\operatorname{Tr}(G_M^{-1}C_M^*C_M).
\]
Then
\[
 1-\pi_M\le\tau_M\le r_M(1-\pi_M^{1/r_M})<r_M,
 \qquad
 \boxed{\mathcal D_k^2\le\tau_M\chi_M
                    \le r_M(1-\pi_M^{1/r_M})\chi_M.}      \tag{CT.26}
\]
To prove every factor, let \(t_1,\ldots,t_{r_M}\in(0,1)\)
be the positive eigenvalues of
\(G_M^{-1/2}\Delta_MG_M^{-1/2}\).  The other eigenvalues
are zero, so
\(\tau_M=\sum_i t_i\) and
\(\pi_M=\prod_i(1-t_i)\).  Induction on the number of
factors proves \(\prod_i(1-t_i)\ge1-\sum_i t_i\).
For positive numbers with a fixed sum, their product is at most the
product with all entries equal: a maximum exists on the closed
nonnegative simplex, lies in its interior when the sum is positive,
and any unequal pair can be replaced by its average to increase
the product, since
\(((a+b)/2)^2-ab=(a-b)^2/4\).  Apply this fact to
\(1-t_i\) to obtain
\(\pi_M\le(1-\tau_M/r_M)^{r_M}\), proving both bounds.
The last two inequalities in (CT.26) now follow from (CT.24).
For \(r_M=1\), \(\tau_M=1-\pi_M\) exactly; for
\(r_M=0\), \(\Delta_M=0\), \(\pi_M=1\), and
\(W_M=0\), so this case requires no exponent \(1/r_M\).
